% =====================================================================
% REVTeX 4-2 manuscript template
%
% Build:    latexmk -pdf main.tex
% Watch:    latexmk -pvc -pdf main.tex
% Clean:    latexmk -c
%
% Bibliography uses biblatex+biber and pulls from the shared master
% library at ../../references/library.bib.
% =====================================================================

\documentclass[
  aps,           % American Physical Society style
  prd,           % Physical Review D (use prl for Phys. Rev. Lett.)
  reprint,       % two-column reprint style; use 'preprint' for single-col drafts
  superscriptaddress,
  nofootinbib,
  amsmath,
  amssymb,
  floatfix,
]{revtex4-2}

% --- Packages -------------------------------------------------------
\usepackage{graphicx}
\usepackage{bm}
\usepackage{physics}      % \dd, \pdv, \dv, \bra, \ket, \expval, ...
\usepackage{tensor}
\usepackage{xcolor}
\usepackage[colorlinks=true,
            linkcolor=blue!60!black,
            citecolor=green!60!black,
            urlcolor=blue!60!black]{hyperref}

% biblatex with biber, APS-physics style
\usepackage[
  backend=biber,
  style=numeric-comp,
  sorting=none,
  giveninits=true,
  uniquename=init,
]{biblatex}
\addbibresource{../../references/library.bib}

% --- Custom macros --------------------------------------------------
\newcommand{\tr}{\operatorname{tr}}
\newcommand{\Tr}{\operatorname{Tr}}
\newcommand{\Sch}{\text{Sch}}                 % Schwarzschild
\newcommand{\GN}{G_{\!N}}                     % Newton's constant
\newcommand{\lp}{\ell_{\mathrm{P}}}           % Planck length
\newcommand{\Mp}{M_{\mathrm{P}}}              % Planck mass

% =====================================================================
\begin{document}

\title{Title of the paper}

\author{Sameer Pant}
\email{spant@clearoneadvantage.com}
\affiliation{Affiliation here}

\date{\today}

\begin{abstract}
One-paragraph abstract. State the problem, the method, the result, and why it
matters --- in that order.
\end{abstract}

\maketitle

% --- Introduction ---------------------------------------------------
\section{Introduction}
\label{sec:intro}

Background and motivation. End the introduction with a roadmap of the paper.

For example, Hawking's original derivation~\cite{Hawking1975} established that
a black hole of mass $M$ emits radiation with a thermal spectrum at temperature
\begin{equation}
  T_H = \frac{\hbar c^3}{8\pi \GN M k_B}.
  \label{eq:hawking-temp}
\end{equation}
The information paradox that follows~\cite{Page1993b, AMPS2013} has been
sharpened in recent years by the island
construction~\cite{Penington2019, AEMM2019, PenningtonShenkerStanfordYang2019}.

% --- Setup ----------------------------------------------------------
\section{Setup}
\label{sec:setup}

State the model, conventions, and notation.

% --- Main calculation ------------------------------------------------
\section{Main calculation}
\label{sec:main}

\subsection{A subsection}

Body.

% --- Results ---------------------------------------------------------
\section{Results}
\label{sec:results}

\begin{figure}[t]
  \centering
  % \includegraphics[width=\columnwidth]{figures/example.pdf}
  \caption{A figure caption that explains what the figure shows and what to
  conclude from it.}
  \label{fig:example}
\end{figure}

% --- Discussion ------------------------------------------------------
\section{Discussion}
\label{sec:discussion}

Limitations, open questions, future directions.

% --- Acknowledgments -------------------------------------------------
\begin{acknowledgments}
We thank \dots\ for useful discussions.
\end{acknowledgments}

% --- Bibliography ----------------------------------------------------
\printbibliography

% --- Appendices ------------------------------------------------------
\appendix

\section{Conventions}
\label{app:conventions}

Metric signature $(-,+,+,+)$. Natural units $\hbar = c = k_B = 1$ unless stated.

\end{document}
